\documentclass{article}

% Four-page, double-blind ML-for-Systems workshop extended abstract.

\usepackage[dblblindworkshop]{neurips_2026}

\workshoptitle{ML for Systems Workshop}

\usepackage[utf8]{inputenc}

\usepackage[T1]{fontenc}

\usepackage{hyperref}

\usepackage{url}

\usepackage{booktabs}

\usepackage{graphicx}

\usepackage{amsmath}

\usepackage{amsfonts}

\usepackage{microtype}

\usepackage{xcolor}

% This is a planning document that compiles in the conference template. Replace

% blue text as evidence becomes available; do not submit with blueprint notes.

\definecolor{blueprint}{RGB}{25,82,145}

\newcommand{\todo}[1]{\textcolor{blueprint}{[TODO: #1]}}

\newcommand{\pilot}[1]{\textcolor{blueprint}{[PILOT: #1]}}

\newcommand{\bytes}[1]{\lvert #1\rvert\_{\mathrm{UTF8}}}

\title{When Worst-Case Prediction Still Compresses: Tokenization and Adversarial Robustness in LLM-Based Lossless Compression}

% Keep the submission anonymous. Add authors only in the camera-ready version.

\author{Anonymous Authors}

\begin{document}

\maketitle

\begin{abstract}

LLM token tables can act as global string compressors: variable-length byte

sequences become fixed-width token IDs that support fast decoding and execution

on encoded data. Predictive LLM compression adds a second layer, entropy-coding

those IDs with next-token probabilities. We show that robustness depends on the

product of both layers---tokens per source byte and coded bits per token---and

cannot be inferred from either token-table compression ratio or perplexity alone.

We construct canonical adversarial sequences and separately attribute

tokenization gain, model surprise, arithmetic-coder loss, dictionary transfer,

and framing. For Qwen2.5-0.5B, full-vocabulary adversarial sequences cost 5.73

times more bits/token than matched text8, but tokenization absorbs enough of that

cost to retain a 51.3\\% ideal byte rate. When restricted to one-byte tokens,

however, the same predictive layer expands data to a 143.4\\% arithmetic payload.

Post-hoc vocabulary masking reduces adversarial model cost by 35.2\\%. These

results identify tokenizer fertility as a robustness boundary and motivate

co-designing token tables, predictive coders, and guarded fallback paths.

\end{abstract}

% -----------------------------------------------------------------------------

% BLUEPRINT AT A GLANCE (delete before submission)

% Target: at most four content pages, including figures and tables. Suggested

% budget: motivation 0.75; formulation and attacks 1.0; evaluation 0.75; results

% 0.9; systems implications, limitations, and conclusion 0.6. Do not assume

% references or supplementary material are excluded from the organizers four-page

% limit; confirm their counting rule before submission.

%

% One-sentence thesis:

% Tokenization is a variable-rate front end to predictive coding, so robust

% evaluation must optimize and report code length per decoded source byte, not

% merely next-token surprisal.

%

% Claim discipline:

% (C1) Formal: token surprisal and source-normalized compression ratio are

%      different objectives whenever decoded token lengths vary.

% (C2) Method: canonical, byte-aware greedy and beam attacks directly optimize

%      entropy-floor or realized arithmetic-code expansion.

% (C3) Empirical: use only after the full multi-model/seed matrix supports it.

%      The current repository establishes a Qwen2.5-0.5B case study, not broad

%      tokenizer or model universality.

% -----------------------------------------------------------------------------

\section{What is doing the compression?}

\`\`Language Modeling Is Compression'' formalizes how next-symbol probabilities

can drive a lossless arithmetic coder and describes tokenization itself as

lossless pre-compression\~\citep{deletang2024language}. This observation implies

that so-called LLM compression combines two distinct mechanisms: a tokenizer

replaces variable-length source strings with token IDs, and a language model

predicts those IDs so that an entropy coder can shorten likely ones. LLMZip is

another instance of this composed design\~\citep{valmeekam2023llmzip}.

Schmidt et al.\~\citep{schmidt2026token} isolate the first mechanism. They use an

LLM vocabulary as a global symbol table in which common byte strings map to

fixed-width IDs and one-byte entries guarantee coverage. Unlike a block-local

code, this stable representation supports equality checks, hashing, joins, and

aggregations before decoding. These properties make token tables particularly

attractive for databases, where strings can remain encoded during common

relational operations and be decoded only when needed. Thus tokenization can be

a useful compression primitive without language-model inference.

What remains unclear is how much of an LLM codec's behavior comes from this

token representation and how much comes from prediction. Delétang et al. compare

ASCII and BPE on natural enwik9, but the two effects have not been isolated under

inputs chosen against realized end-to-end size. This matters because prediction

loss is measured per token while compression is measured per source byte: a poor

probability model may still sit inside a pipeline that compresses when its tokens

represent many bytes, whereas low-fertility tokenization can expose expansion.

We therefore study the tokenizer--predictor--coder pipeline as a two-layer system.

\paragraph{Contributions.}

We (i) factor end-to-end rate into tokenizer fertility and predictive coding

cost; (ii) construct lossless attacks against token-, byte-, and realized-code

objectives; and (iii) use matched natural and adversarial inputs to show when the

token layer absorbs predictive surprise and when it exposes expansion.

\paragraph{Scope.}

We hold one pretrained tokenizer/model fixed and vary canonical inputs

adversarially. Our evidence is a Qwen2.5-0.5B case study and does not contradict

the token-table paper's corpus-level results.

\section{Compression objectives and attacks}
\label{sec:attacks}

Our central distinction is between prediction difficulty and end-to-end compression
failure. A token can be extremely unlikely under the language model while decoding
to many source bytes, so maximizing token surprisal need not maximize compressed
size per input byte. We therefore evaluate three increasingly end-to-end attack
objectives side by side.

Let $h$ denote the current token history and let $v$ be a candidate next token. The
first objective ignores token length and selects the least probable valid extension:
\begin{equation}
  \operatorname{MinProb}(v;h)=-\log_2 p_\theta(v\mid h).
  \label{eq:minprob}
\end{equation}
This is the conventional token-level worst case. It maximizes ideal coded bits per
token, but not necessarily bits per source byte.

To account for variable-length tokenization, define the incremental decoded byte
length
\begin{equation}
  \Delta b(v;h)=
  \bytes{\operatorname{decode}(h\mathbin\Vert v)}-
  \bytes{\operatorname{decode}(h)}.
\end{equation}
The byte-aware greedy objective is
\begin{equation}
  \operatorname{SurprisalPerByte}(v;h)=
  \frac{-\log_2 p_\theta(v\mid h)}{\Delta b(v;h)}.
  \label{eq:spb}
\end{equation}
It directly targets the ideal entropy-coding cost per newly decoded source byte.

Finally, for a complete canonical token sequence $x$, let
$\operatorname{SerializeCodec}(x)$ denote the actual serialized codec output. The
end-to-end objective is
\begin{equation}
  \operatorname{ActualRate}(x)=
  \frac{\left|\operatorname{SerializeCodec}(x)\right|}
       {8\bytes{\operatorname{decode}(x)}}.
  \label{eq:actualrate}
\end{equation}
A realized-size beam search approximately maximizes Eq.~\ref{eq:actualrate} by
cloning coder state for candidate continuations. This objective captures arithmetic
coding and serialization effects that ideal NLL does not.

\subsection{Why the objectives differ}
For a sequence of $N$ tokens decoding to $U$ source bytes, define fertility
$f(x)=U/N$. If
\begin{equation}
  L(x)=-\sum_{t=2}^{N}\log_2 p_\theta(x_t\mid x_{<t})
\end{equation}
is the sequence NLL, then the ideal predictive rate is
\begin{equation}
  R_{\rm ideal}(x)=\frac{L(x)}{8U}
  \approx
  \underbrace{\frac{L(x)}{N}}_{\text{bits/token}}
  \underbrace{\frac{N}{U}}_{\text{tokens/byte}}
  \frac{1}{8}.
  \label{eq:factorization}
\end{equation}
Thus high bits/token can be offset by high bytes/token. This is the robustness
mechanism we test: tokenization can buffer severe prediction failure when
adversarial tokens represent many source bytes.

\subsection{Threat model and generation policies}
The adversary is white-box and has access to the tokenizer, model logits, and coding
algorithm. It chooses a length-$N$ token sequence after a fixed seed, subject to a
candidate alphabet $G$. Because arbitrary token IDs need not survive decoding and
retokenization, every accepted prefix must satisfy
\begin{equation}
  T(\operatorname{decode}(x_{1:t}))=x_{1:t}.
  \label{eq:canonical}
\end{equation}
Final decoding must recover exactly the same IDs and bytes.

We compare five generation policies. \emph{Natural text} is the non-adversarial
corpus baseline. \emph{Random canonical} uniformly samples valid canonical
extensions and controls for unusual token sequences without optimizing against the
model. \emph{Minimum probability} greedily maximizes Eq.~\ref{eq:minprob}.
\emph{Maximum surprisal/byte} greedily maximizes Eq.~\ref{eq:spb}.
\emph{Realized-size beam} searches over multiple canonical continuations to
approximately maximize Eq.~\ref{eq:actualrate}.

The \emph{one-byte constraint} is not a separate attack objective. It is an ablation
of the candidate alphabet in which $G$ contains only canonical tokens contributing
one decoded byte. This fixes fertility near one byte/token and removes the
multi-byte-token buffering mechanism. Under an exact one-byte constraint,
MinProb and SurprisalPerByte induce the same greedy ranking.

\paragraph{Global mask.}
When evaluating occurring-token dictionaries, we construct each adversarial
sequence once under the full distribution and replay identical IDs and histories
under the post-hoc occurring-token mask. This isolates probability renormalization;
it is not a mask-adaptive attack.

\section{Experimental design}
\label{sec:design}

\paragraph{Model and codec.}
We evaluate Qwen2.5-0.5B with its 151,643-token vocabulary. The natural-text control
uses the first 100,000 text8 tokens with context length 1,000, retained context 100,
KV caching, and arithmetic coding. It uses 16 independently seeded batches, so
99,984 tokens are predicted. Model weights are treated as shared decoder state.

\paragraph{Adversarial runs.}
The existing full-vocabulary experiment contains three 1,000-token minimum-
probability sequences initialized with token IDs 785, 32, and 641. All completed
sequences satisfy Eq.~\ref{eq:canonical} and decode to 6,988, 6,812, and 6,943 UTF-8
bytes. Their post-hoc occurring-token dictionaries contain 592, 606, and 587 tokens.
We report adversarial means with sample standard deviation across the three starts.
The matched text8 run is deterministic for each dictionary policy and therefore has
no error bar.

The decisive evaluation runs the generation policies side by side under the same
sequence-length, seed, canonicality, context, and codec settings. This controls the
search objective while keeping the compression pipeline fixed. The one-byte
experiment repeats the relevant attacks with the restricted candidate alphabet to
isolate the contribution of token fertility.

\paragraph{Reported quantities.}
For every input or attack we report the same decomposition:
\begin{align}
  \text{Bits/token} &= \frac{L(x)}{N_{\rm pred}},\\
  \text{Bytes/token} &= \frac{\bytes{\operatorname{decode}(x)}}{N},\\
  \text{Ideal bits/byte} &= \frac{L(x)}{\bytes{\operatorname{decode}(x)}},\\
  \text{Actual bits/byte} &=
  \frac{\left|\operatorname{SerializeCodec}(x)\right|}
       {\bytes{\operatorname{decode}(x)}},\\
  \text{Final ratio} &= \frac{\text{Actual bits/byte}}{8}.
\end{align}
A final ratio below one denotes compression relative to raw UTF-8; a value above
one denotes expansion. Where a realized codec run has not yet been measured, we
report ``--'' rather than substituting an ideal rate.

\begin{table*}[t]
  \centering
  \small
  \caption{Side-by-side attack evaluation. Bits/token measures model prediction
  difficulty; bytes/token measures tokenizer fertility; ideal bits/byte combines
  the two; actual bits/byte includes the realized codec; and the final ratio is
  actual size divided by the original 8-bit UTF-8 size. The one-byte row is a
  constrained-alphabet ablation rather than a distinct objective.}
  \label{tab:decisive}
  \begin{tabular}{@{}lrrrrr@{}}
    \toprule
    Attack / input & Bits/token & Bytes/token & Ideal bits/byte & Actual bits/byte & Final ratio \\
    \midrule
    Natural text & 4.954 & \todo{measure} & \todo{derive} & \todo{measure} & 0.14405 \\
    Random canonical & \todo{run} & \todo{run} & \todo{run} & \todo{run} & \todo{run} \\
    Minimum probability & $28.409\pm0.077$ & 6.914 & 4.108 & \todo{run} & \todo{run} \\
    Maximum surprisal/byte & \todo{run} & \todo{run} & \todo{run} & \todo{run} & \todo{run} \\
    Realized-size beam & \todo{run} & \todo{run} & \todo{run} & \todo{run} & \todo{run} \\
    \midrule
    One-byte constrained & \todo{report matched attack} & 1.000 & \todo{report} & \todo{report} & \todo{report} \\
    \bottomrule
  \end{tabular}
\end{table*}

\section{Results}
\label{sec:results}

\paragraph{Prediction failure does not imply compression failure.}
The existing minimum-probability full-vocabulary attack reaches
$28.409\pm0.077$ ideal bits per predicted token, $5.734\times$ the matched text8
value of 4.954 bits/token. Because a fixed ID for the 151,643-token vocabulary
requires 18 bits, prediction makes this adversarial token stream substantially more
expensive than fixed-width IDs. Yet the sequences decode to an average 6.914 source
bytes/token. Dividing model cost by this fertility gives approximately 4.108 ideal
bits/source byte, or 51.346\% of raw UTF-8. The model therefore fails badly at the
token level while the composed tokenizer--predictor system still has an ideal
compression rate below one.

\paragraph{Byte-aware attacks test the relevant compression objective.}
Minimum probability optimizes only the numerator of Eq.~\ref{eq:factorization}.
Maximum surprisal/byte additionally penalizes tokens that obtain high surprisal only
by representing many bytes, while realized-size beam targets the serialized output
itself. Table~\ref{tab:decisive} is therefore the primary comparison: it determines
whether progressively more end-to-end objectives increase ideal or realized
bits/source byte. The corresponding full-vocabulary SurprisalPerByte and realized-
size beam measurements are required before making a claim that MinProb is the
strongest compression attack.

\paragraph{Removing tokenization gain exposes predictive expansion.}
The existing one-byte-ASCII ablation fixes fertility at one byte/token. In this
setting greedy minimum probability and surprisal-per-byte coincide. The measured
model floor is 135.48\% of raw bytes and the arithmetic payload is 143.36\%; an NLL
beam reaches 136.70\% ideal and 144.14\% arithmetic payload. Thus the predictive
layer cannot beat the original 8-bit bytes once the tokenizer is prevented from
absorbing model surprise through multi-byte tokens. The narrow 32-token beam
validates this mechanism and the accounting, not global optimality.

\paragraph{Dictionary masking changes prediction cost, not the attack sequence.}
Post-hoc occurring-token dictionaries contain 592, 606, and 587 symbols for the
three existing adversarial sequences. Replaying identical IDs and histories under
these masks lowers adversarial cost from $28.409\pm0.077$ to
$18.402\pm0.833$ bits/token, a 35.2\% reduction, and lowers the byte-normalized
ideal rate from 51.346\% to 33.284\%. It remains $3.806\times$ the matched masked-
text8 cost of 4.835 bits/token. Because generation is unchanged, this experiment
isolates renormalization rather than adaptive robustness.

\paragraph{Realized overhead must be reported separately from the ideal floor.}
For text8, model-induced floors are 11.631\% of raw UTF-8 with the full vocabulary
and 11.352\% with occurring tokens. Finite coding and metadata raise realized totals
to 14.405\% and 13.434\%, corresponding to $6.942\times$ and $7.444\times$
compression. The occurring-token bitmap costs 10,818 bytes (2.03\% of raw size),
less than its coding gain at 100,000 tokens. These gaps motivate keeping ideal and
actual columns separate in Table~\ref{tab:decisive}.

\section{Systems implications and limitations}

\label{sec\:limitations}

The systems lesson is to preserve token IDs as the global, queryable in-memory

representation and treat predictive arithmetic coding as an optional block-level

storage or transport layer. Source-normalized accounting identifies whether

fragility comes from the tokenizer, predictor, quantizer, dictionary, or

container. For untrusted blocks, a guard can store the smaller of predictive and

conventional encodings plus a selector bit, bounding size by the fallback plus one

bit; its decoding granularity, latency, and throughput remain to be measured.

Our evidence is limited to one model, synthetic white-box inputs, three starts,

finite context, approximate search, and one coder/container. It does not establish

global worst cases or generalize to multilingual text or code. The attack could

enable storage or compute amplification; practical mitigations include size caps,

early-abort expansion detection, bounded inference, and fallback codecs.

\section{Conclusion}

Robust learned compression is an end-to-end systems problem. Canonical byte-aware attacks and explicit cost attribution show that the tested

adversarial sequences are 3.81--5.73 times harder than matched text8, while

guarded fallback offers a principled path from diagnosis to robust operation.

\begin{ack}

\todo{Camera-ready only: funding and competing-interest disclosures. The supplied

style hides this environment in anonymous submission mode.}

\end{ack}

{\small

\begin{thebibliography}{5}

\bibitem[Delétang et al.(2024)]{deletang2024language}

Grégoire Delétang, Anian Ruoss, Paul-Ambroise Duquenne, Elliot Catt, Tim

Genewein, Christopher Mattern, Jordi Grau-Moya, Li Kevin Wenliang, Matthew

Aitchison, Laurent Orseau, Marcus Hutter, and Joel Veness.

\newblock Language modeling is compression.

\newblock In \emph{*ICLR*}, 2024.

\bibitem[Schmidt et al.(2026)]{schmidt2026token}

Tobias Schmidt, Nicolas Schmitt, Thomas Neumann, and Andreas Kipf.

\newblock Accelerating string-heavy queries with LLM token tables.

\newblock Manuscript, 2026.

\bibitem[Shannon(1948)]{shannon1948}

Claude E. Shannon.

\newblock A mathematical theory of communication.

\newblock \emph{*Bell System Technical Journal*}, 27:379--423, 623--656, 1948.

\bibitem[Valmeekam et al.(2023)]{valmeekam2023llmzip}

Chandra S. K. Valmeekam, Krishna Narayanan, Dileep Kalathil,

Jean-François Chamberland, and Srinivas Shakkottai.

\newblock LLMZip: Lossless text compression using large language models.

\newblock arXiv:2306.04050, 2023.

\bibitem[Witten et al.(1987)]{witten1987}

Ian H. Witten, Radford M. Neal, and John G. Cleary.

\newblock Arithmetic coding for data compression.

\newblock \emph{*Communications of the ACM*}, 30(6):520--540, 1987.

\end{thebibliography}

}

% Supplementary planning notes (do not count toward the four-page narrative):

% - Prove the token/byte objective mismatch and give attack pseudocode.

% - Record configurations, revisions, commands, seeds, hardware, compute, and

%   licenses; report full per-run tables and search-budget ablations.

% - Document exact seed-token, quantization, bitmap, payload, and header costs.

% - Repository map: src/adversarial.py; src/compression\_attacks.py;

%   scripts/run\_compression\_attacks.py; evaluation/notebooks/neurips/

%   adversarial\_worst\_case.ipynb; and artifacts/runs/adversarial/

%   qwen\_05b\_ascii\_bytes\_n1000.

% - Confirm whether references and the NeurIPS checklist are outside the

%   organizers four-page limit. Keep the checklist unless organizers explicitly

%   say their workshop does not require it.

\newpage

\input{checklist.tex}

\end{document}